\begin{center}
	\large{\textbf{GRAVITASI EINSTEIN DAN EINSTEIN-MAXWELL PADA 2+1 DIMENSI DALAM TEORI GAUGE CHERN-SIMONS}}
\end{center}

\vspace{1em}

\noindent
\begin{tabular}{@{}lcl@{}}
	Nama       & : & Agung Sedayu Septiawan \\
	NRP        & : & 5001221075 \\
	Departemen & : & Fisika \\
	Pembimbing & : & Dr. rer. nat. Bintoro Anang Subagyo \\
\end{tabular}

\vspace{1em}

\noindent\textbf{Abstrak}

\par Penelitian ini merumuskan gravitasi Einstein di 2+1 dimensi sebagai teori gauge Chern-Simons dengan grup gauge $\mathrm{SO}(2,2)$ dan menerapkannya pada lubang hitam BTZ serta perluasan Einstein-Maxwell. Melalui ansatz Witten, koneksi spin dan dreibein digabung menjadi koneksi gauge tunggal bernilai pada aljabar Lie $\mathfrak{so}(2,2)$ yang dikonstruksi dari prinsip pertama; variasi aksi Chern-Simons mereproduksi secara persis kondisi bebas torsi dan persamaan medan Einstein dengan konstanta kosmologi negatif, membuktikan ekuivalensi kedua formulasi. Solusi lubang hitam BTZ beserta horizon peristiwanya diperoleh secara analitik dalam formalisme dreibein, dengan kasus tanpa rotasi dianalisis sebagai limit yang mereduksi metrik ke bentuk diagonal. Perluasan ke sektor Einstein-Maxwell dilakukan dengan menambahkan aksi Maxwell sebagai sumber materi; persamaan Maxwell diselesaikan secara umum dan tensor energi-momentum elektromagnetik dihitung eksplisit dalam notasi dreibein. Analisis konsistensi menunjukkan bahwa kesamaan struktural geometris yang dipaksakan oleh ansatz BTZ berbenturan dengan asimetri tensor energi-momentum elektromagnetik, sehingga solusi yang sekaligus bermuatan dan berotasi tidak ada dalam teori Einstein-Maxwell murni di 2+1 dimensi.

\vspace{1em}

\noindent\textbf{Kata kunci}: \textit{Gravitasi 2+1 dimensi, Chern-Simons, BTZ, Einstein-Maxwell}

\newpage
\thispagestyle{empty}
\vspace*{\fill}
    \begin{center}
	Halaman ini sengaja dikosongkan
    \end{center}
\vspace*{\fill}
\pagebreak
